\section{Introduction}
\label{sec:intro}

Mechanistic interpretability has a mature toolbox for \emph{validating}
hypotheses about model internals: activation patching quantifies where
information is causally used \citep{vig2020,meng2022,zhang2024}, and
distributed alignment search (DAS) tests whether a proposed causal variable
is implemented in a linear subspace \citep{geiger2024das,wu2023boundless}.
What these methods do not provide is the hypothesis itself. In practice the
analyst chooses a concept---gender, sentiment, a high-level algorithm---and
constructs contrastive data by hand, with no independent guarantee that the
chosen concept is the one the task actually requires
\citep{hase2023,lanham2023}. Validation is rigorous; generation is
intuition.

We study a setting where generation can be made independent and verifiable.
Rosetta Stone puzzles from the UK Linguistics Olympiad (UKLO) present a
small set of translation pairs in an unfamiliar language---for example,
Gilbertese \emph{E nooria te uaa te aomata} `The man saw the fruit'---and
ask for translations of held-out sentences \citep{bozhanov2013,lingoly}.
The puzzles are self-contained: every rule needed to answer is inferable
from the given pairs. This yields built-in verification. We use an Answer
Set Programming (ASP) solver \citep{gelfond1988,gebser2014} to derive a
minimal set of grammatical rules that provably regenerates every context
pair and uniquely translates the target question. Each \emph{certified}
rule (an object-agreement suffix, a tense particle, a noun-class concord)
is a discrete, machine-verified hypothesis about what a model solving the
puzzle must compute---selected by the task structure, not by the analyst,
and converted automatically into counterfactual minimal pairs for causal
testing.

We apply this framework to \texttt{Llama-3.1-8B} on four questions from
three typologically distinct languages (Gilbertese, Ndebele, Tshiluba),
yielding 29 certified rules. Our contributions:
\begin{enumerate}[leftmargin=*,itemsep=0pt,topsep=2pt]
  \item A framework for \textbf{certified hypothesis generation}: causal
  variables derived from task structure and verified by a solver,
  decoupling hypothesis generation from both researcher intuition and the
  model under study (\S\ref{sec:generation}).
  \item A quantitative characterization of the generation pipeline itself:
  certification rates, repair-loop dynamics, and uniqueness of the
  certified grammars. \TODO{pending extraction experiments; see
  \S\ref{sec:results-generation}}
  \item \textbf{Causal validation} using standard interventions
  (\S\ref{sec:validation}): rule implementations stratify by
  depth---morphosyntax early (L0--L4), lemma selection near the output
  (L20--L31); individual rules admit causally sufficient linear subspaces
  of rank $\le 8$; interventions on independent rules commute (100\% of 16
  rule pairs); and a dissociation---every rule is linearly decodable at the
  final token, where interventions nonetheless have near-zero causal
  effect (\S\ref{sec:results-validation}).
\end{enumerate}

We use only off-the-shelf validation machinery and claim no novelty there;
the contribution is the provenance and verifiability of the hypotheses, and
the resulting evidence about how in-context grammatical rules are
implemented.
